% TOP_PROBLEM: {"id": 1486, "title": "Whitehead Conjecture", "category_id": 7, "category": {"id": 7, "name": "topology", "display_name": "Topology", "description": "Properties preserved under continuous deformations.", "slug": "topology", "order_index": 7, "created_at": "2026-07-31T15:26:25.671Z"}, "status": "open"}
% FIRST_POSTED: 2026-09-27
% !TeX program = lualatex
% ATTEMPT_STATUS: unresolved
\documentclass[11pt]{article}
\usepackage[margin=25mm]{geometry}
\usepackage{fontspec}
\setmainfont{Latin Modern Roman}
\usepackage{amsmath,amssymb,mathtools}
\usepackage{unicode-math}
\setmathfont{Latin Modern Math}
\usepackage{xurl,hyperref}
\hypersetup{colorlinks=true,urlcolor=blue}
\setcounter{secnumdepth}{0}
\setlength{\parindent}{0pt}
\setlength{\parskip}{5pt}
\setlength{\emergencystretch}{3em}
\begin{document}
\section*{\texorpdfstring{Whitehead Conjecture}{Whitehead Conjecture}}\label{1486}
\subsection{Definitions and mathematical statement}
A connected CW-complex $X$ is aspherical if $\pi_i(X)=0$ for all $i\ge2$, equivalently its universal cover is contractible. A subcomplex is a union of cells closed under taking attaching boundaries. Whitehead's conjecture asks
\[
 X\text{ aspherical},\quad\dim X\le2,\quad
 Y\subseteq X\text{ a connected subcomplex}
 \quad\Longrightarrow\quad Y\text{ aspherical}.
\]
Neither finiteness of the complexes nor a special fundamental group is assumed in the recorded scope. In presentation language, a presentation complex has one two-cell per relator; removing relators gives the corresponding subpresentation version. Merely knowing that the inclusion induces an injection on fundamental groups would not by itself be a definition of asphericity.
\par\smallskip\textit{Formulation and concept references:} \href{https://www.mathnet.ru/php/archive.phtml?jrnid=im&option_lang=rus&paperid=9597&wshow=paper}{[S1]}.

\subsection{Short English statement}
Must every connected subcomplex of an aspherical two-dimensional cell complex also have contractible universal cover?
\subsection{Research attempt}
\paragraph{Scope and proof-claim audit (27 September 2026).}
The assertion concerns arbitrary connected subcomplexes, including
infinite ones. Asphericity is a statement about the universal cover,
not just ordinary second homology or the Euler characteristic of the
original complex. The dimension bound two will be used explicitly.

Mikhovich's 2025 primary text [S1] concerns rational and $p$-adic analogues.
Its completed or prounipotent conclusions are not silently substituted
for the integral discrete statement. There are also primary manuscripts
claiming affirmative answers to the exact classical question: Pasku,
arXiv:2107.12293v1 (26 July 2021), and Kawauchi,
arXiv:2303.04368v2 (9 April 2024) [E1,E2]. Their current abstract/version
pages were inspected and the claims are recorded as such. This attempt
does not independently certify their full arguments or infer consensus
acceptance from the presence of an abstract. The inherited disputed-claim
metadata is preserved, and our own status remains unresolved.

The approach below proves the conclusion when the inclusion is injective
on fundamental groups, when the subcomplex has a retraction, and for
subcomplexes of finite collapsible two-dimensional simplicial complexes.
It then isolates the precise covering-space mistake in a tempting proof
of the full conjecture. No finiteness assumption is inserted into the
general target.

\paragraph{A useful two-dimensional asphericity criterion.}
\textbf{Lemma 397.1.} A connected two-dimensional CW-complex $Z$ is
aspherical if and only if its universal cover $\widetilde Z$ has
$H_2(\widetilde Z;\mathbb Z)=0$.

\emph{Proof.} The forward direction follows from contractibility.
Conversely the cover is simply connected, so $H_1(\widetilde Z)=0$.
Cellular homology gives zero homology in all degrees above two because
there are no higher-dimensional cells. By hypothesis its reduced
homology therefore vanishes in every degree.

If any homotopy group of $\widetilde Z$ were nonzero, take the least
positive index $j$ for which this happens. It is at least two. The
Hurewicz theorem for a $(j-1)$-connected CW-complex identifies that
first nonzero homotopy group with $H_j$, a contradiction. Thus all
homotopy groups vanish. The map to a point is a weak homotopy equivalence
of CW-complexes and is a homotopy equivalence by Whitehead's theorem.
The cover is contractible.\hfill$\square$

Here Whitehead's general weak-equivalence theorem is an established
input, distinct from the open asphericity conjecture. This proof also
explains why checking $\pi_2$ is enough for a two-dimensional complex:
Hurewicz identifies $\pi_2(\widetilde Z)$ with $H_2(\widetilde Z)$,
and then the same argument eliminates all higher groups.

\paragraph{Subcomplexes inherit ordinary second-homology vanishing upstairs.}
Let $E$ be a contractible CW-complex of dimension at most two and
$Z\subseteq E$ a subcomplex. In cellular chains there is no group
$C_3(E)$, so
\[
 H_2(E)=\ker(\partial_2^E)=0.
\]
The map $C_2(Z)\to C_2(E)$ is injective, being the inclusion of the
free abelian group on a subset of the cells. Boundaries commute with
inclusion. Therefore every two-cycle of $Z$ is a two-cycle of $E$,
and must be zero. Hence
\[
 H_2(Z;\mathbb Z)=0.
 \tag{397.1}
\]
The same chain argument works with any coefficient group. It does
not yet say that $H_2(\widetilde Z)=0$: the universal cover of $Z$
need not embed in $E$.

There is additional integral information when $Z$ is connected. The
long exact sequence of $(E,Z)$ contains
\[
 0=H_2(E)\longrightarrow H_2(E,Z)
 \longrightarrow H_1(Z)\longrightarrow H_1(E)=0.
\]
Thus $H_1(Z)\cong H_2(E,Z)$. Because the relative cellular chain
complex also has no degree-three group, $H_2(E,Z)$ is a subgroup of
the free abelian group $C_2(E,Z)$. Subgroups of free abelian groups
are free abelian, including in infinite rank. Consequently
\[
 H_1(Z;\mathbb Z)\text{ is free abelian}.
 \tag{397.2}
\]
This is a necessary condition on subcomplexes of contractible
\emph{two}-complexes; it is not a claim that their fundamental groups
are free. Abelianization loses commutator information.

For example the projective plane cannot occur as such a subcomplex:
its first integral homology is $\mathbb Z/2$, contrary to (397.2).
The second homology test alone would miss it, since
$H_2(\mathbb{RP}^2;\mathbb Z)=0$ although its universal cover is $S^2$.
This also supplies a concrete warning against replacing universal-cover
homology by ordinary integral homology in Lemma 397.1.

\paragraph{A complete subcase: injectivity on fundamental groups.}
\textbf{Proposition 397.2.} If $Y\subseteq X$ is as in the problem
and $\pi_1(Y)\to\pi_1(X)$ is injective, then $Y$ is aspherical.

\emph{Proof.} Let $p:\widetilde X\to X$ be the universal cover,
and choose a component $Z$ of $p^{-1}(Y)$. It is a subcomplex of the
contractible two-complex $\widetilde X$. The restriction $Z\to Y$
is the connected covering corresponding to the subgroup
\[
 \ker\bigl(\pi_1(Y)\to\pi_1(X)\bigr).
\]
Indeed a loop in $Y$ lifts to a closed loop in $\widetilde X$ exactly
when its image in $\pi_1(X)$ is trivial. Injectivity makes this
subgroup trivial, so $Z$ is the universal cover of $Y$.
Equation (397.1) gives $H_2(Z)=0$, and Lemma 397.1 finishes the proof.
\hfill$\square$

The proposition uses more than the definition of asphericity: it is
a proof that, in this dimension and ambient setting, injectivity is
a sufficient additional hypothesis. The statement in the catalogue
correctly warns that injectivity by itself is not the definition.

It is not a necessary hypothesis either. Take $X$ to be a closed disc
with its boundary circle as a subcomplex $Y$. The disc is contractible,
and the circle is aspherical, but the induced map
$\mathbb Z\to0$ is not injective. A proof of the general conjecture
cannot simply try to show that every inclusion is injective; this
smallest example already disproves that intermediate assertion.

\paragraph{Retractions give a second complete mechanism.}
If there is a continuous retraction $r:X\to Y$ with $r|_Y=\mathrm{id}$,
then for every $j\ge1$ the composition
$\pi_j(Y)\to\pi_j(X)\to\pi_j(Y)$ is the identity, using a base point
in $Y$. If $X$ is aspherical, this forces $\pi_j(Y)=0$ for all
$j\ge2$. Thus retracts of aspherical spaces are aspherical, without
any dimensional restriction. Alternatively the $j=1$ statement supplies
the injectivity hypothesis of Proposition 397.2 in the present dimension.

For a presentation complex, a proposed retraction must respect every
attaching map of a removed cell. Sending extra edges to paths in $Y$
only gives a map on the one-skeleton. It extends over each two-cell
only when the corresponding boundary loop becomes nullhomotopic in $Y$.
Removing a relator does not automatically make that relator trivial
in the smaller presentation. Assuming such an extension would hide
the main relation issue rather than solve it.

\paragraph{A constructive class: subcomplexes of collapsible two-complexes.}
\textbf{Proposition 397.3.} If $X$ is a finite collapsible simplicial
complex of dimension at most two, every connected subcomplex $Y$
collapses to a graph and is aspherical.

\emph{Proof.} Fix a collapse sequence for $X$. Record just the pairs
consisting of a triangle and its free edge, in their original order.
Executing these pairs while omitting all edge--vertex collapses remains
valid: an omitted edge--vertex collapse removes no triangle, and cannot
be needed to make an edge a face of only its paired triangle. Thus this
sequence removes every triangle from $X$.

Apply the same ordered list to $Y$, performing a pair exactly when its
triangle is still in $Y$. Its paired edge then also lies in $Y$ because
$Y$ is a subcomplex. No other remaining triangle of $Y$ contains that
edge: such a triangle would also have been present in the ambient
sequence at that stage and would contradict freeness there. Hence the
pair is a valid elementary collapse of $Y$. If the triangle is absent,
do nothing, even if its paired edge happens to belong to $Y$.

At the end every triangle initially in $Y$ has been removed. The
remaining subcomplex is a finite graph. Elementary collapses are homotopy
equivalences, so $Y$ has the homotopy type of this graph.
A connected graph has a tree as its universal cover; a tree is
contractible. Thus $Y$ is aspherical.\hfill$\square$

The result is deliberately about a supplied collapse sequence. Contractible
and collapsible are different conditions; no theorem here converts an
arbitrary contractible two-complex into a collapsible one while preserving
its subcomplexes. Even replacing $X$ by a homotopy-equivalent graph or
point would not preserve the inclusion diagram needed for this proof.

The same triangle-removal argument works whenever there is a sequence
of free-edge collapses eliminating all two-simplices, even if the final
graph has cycles and the ambient complex is not contractible. This gives
a checkable sufficient certificate of asphericity for every subcomplex
in that class. Finding no such collapse does not certify nonasphericity.

\paragraph{The exact obstruction in the universal-cover attack.}
In the general problem, again take a component $Z$ of $p^{-1}(Y)$.
Equations (397.1) and (397.2) still hold, and
\[
 \pi_1(Z)\cong K,
 \qquad K=\ker\bigl(\pi_1(Y)\to\pi_1(X)\bigr).
\]
If $K$ is nontrivial, $Z$ is not the universal cover of $Y$.
The universal cover $\widetilde Y\to Z$ gives a lifted map
$\widetilde Y\to\widetilde X$, but it need not be an embedding.
Different lifts of cells can map onto the same cell of $\widetilde X$.
Consequently the resulting map on free cellular chain groups need not
be injective. The step ``a two-cycle upstairs is a two-cycle in the
contractible ambient cover, hence zero'' fails exactly at this point.

This can already be seen at dimension one in the disc-boundary example.
The preimage of the circle in the universal cover of the disc is the
circle itself. Its universal cover is a line. The lifted map from that
line into the disc repeatedly traverses the same boundary; it is not an
inclusion of a line-shaped subcomplex. The example does not violate
asphericity, but it visibly destroys the asserted embedding on which
the overly short proof would rely.

One must instead prove that the universal cover of $Z$ has no two-cycles,
or obtain another argument killing its second homotopy. Ordinary
homology of $Z$ alone cannot do that. In particular (397.2) only says
$K/[K,K]$ is free abelian; it supplies no presentation of $K$ as a free
group and no reason that $Z$ is a $K(K,1)$.

\paragraph{Why the dimension bound is indispensable in the chain proof.}
The boundary sphere $S^2\subset D^3$ is a subcomplex of a contractible
three-dimensional CW-complex and is not aspherical. In cellular chains
its nonzero two-cycle is killed by the boundary of the three-cell.
Thus in dimension three, $H_2(X)=0$ no longer implies injectivity of
$\partial_2^X$: it only says that its kernel is the image of
$\partial_3^X$. The argument proving (397.1) relied on $C_3(X)=0$.

This example also has injective fundamental-group map, since both
fundamental groups are trivial. It shows why Proposition 397.2 cannot
be promoted to arbitrary-dimensional aspherical ambient complexes.
The combination of dimension two, cellular inclusion, and injectivity
is what makes that proposition work.

\paragraph{Finite witnesses do not justify a finite ambient reduction.}
Every map from a sphere into a CW-complex has compact image, and every
compact subset of a CW-complex is contained in a finite subcomplex.
Consequently, if every finite connected subcomplex of a given complex
$Y$ is aspherical, then $Y$ is aspherical: a sphere map lies in one
such finite subcomplex and is nullhomotopic there. The same argument
works for all indices $j\ge2$.

This observation reduces a conclusion about one fixed $Y$ to its finite
subcomplexes, but it does not manufacture finite aspherical ambient
subcomplexes of an infinite aspherical $X$. A finite set of cells in
$X$ can support spheres that are killed only after adding other cells;
there is no argument here that the process of adding all needed
nullhomotopies terminates after finitely many additions. Likewise,
a compact nullhomotopy for one sphere gives only one finite enlargement,
not a finite enlargement resolving every sphere simultaneously.

The distinction is the same quantifier issue in another form: for each
sphere there exists a finite witness is weaker than existence of one
finite aspherical complex accommodating all spheres. We use neither
an unproved compactness theorem of that strength nor a finiteness
hypothesis absent from the record.

\paragraph{Presentation chains and what integer matrices omit.}
For a finite presentation complex, ordinary cellular $\partial_2$ is
the matrix of exponent sums of generators in relators. It computes
ordinary second homology, but the universal-cover boundary matrix has
coefficients in $\mathbb Z[\pi_1(Y)]$, with coefficients recording
where lifted edge traversals occur. Augmenting those group-ring
coefficients to integers forgets exactly that location information.

For the projective-plane presentation $\langle a\mid a^2\rangle$,
the ordinary matrix is $(2)$ and has zero integer kernel. Upstairs,
with $a^2=1$, the boundary coefficient is $1+a$, and
\[
 (1-a)(1+a)=1-a^2=0.
\]
The nonzero element $1-a$ produces a two-cycle. Geometrically the
universal cover is $S^2$, so this calculation is consistent with its
nonzero second homology. It is a diagnostic for the matrix shortcut,
not a candidate subcomplex of the ambient class; (397.2) already
excludes that candidate.

In the torus presentation $\langle a,b\mid aba^{-1}b^{-1}\rangle$,
the ordinary two-boundary is zero although the torus is aspherical.
Its universal cover is the square tiling of the plane and is contractible.
Thus ordinary second homology can be nonzero in a genuinely aspherical
complex, as well as zero in a nonaspherical one. The correct matrix is
the one upstairs, and a rank computation after augmentation is no
substitute for proving its injectivity.

\paragraph{Verification and first unresolved step.}
The finite diagnostic checks free-edge collapse certificates and the
restriction procedure on subcomplexes of a triangulated disc. It also
checks the group-ring identity in the cyclic order-two example and
the distinct ordinary boundary matrices. These calculations verify
specific certificates and expose false implications; they do not decide
asphericity from finitely many homology ranks.

\textbf{Outcome: UNRESOLVED.} The injective, retract and collapsible
subcases are proved, and the ambient-cover argument gives explicit
homological restrictions without finiteness assumptions. The missing
general step is to control the nontrivial kernel cover $Z$ strongly
enough to show $H_2(\widetilde Z)=0$. None of the chain inclusions,
ordinary homology tests, or unverified external affirmative claims is
used to assert that this step has been completed.

\subsection{Sources}
\begin{itemize}
\item[S1]A. M. Mikhovich, Rational and p-adic analogues of J. H. C. Whitehead's conjecture, Izvestiya: Mathematics 89:2 (2025), section 1.
\url{https://www.mathnet.ru/php/archive.phtml?jrnid=im&option_lang=rus&paperid=9597&wshow=paper}.
\textit{Formulation source.}
\item[E1]Elton Pasku, \emph{An answer to the Whitehead asphericity
question}, arXiv:2107.12293v1, 26 July 2021.
\url{https://arxiv.org/abs/2107.12293}.
\textit{Primary abstract and version history inspected 27 September 2026.
Its affirmative claim is recorded, not independently certified here.}
\item[E2]Akio Kawauchi, \emph{Whitehead aspherical conjecture via ribbon
sphere-link}, arXiv:2303.04368v2, 9 April 2024.
\url{https://arxiv.org/abs/2303.04368}.
\textit{Primary abstract and version history inspected 27 September 2026.
The claim covers arbitrary aspherical two-complexes; the complete proof
is not independently validated in this attempt.}
\item[E3]A. M. Mikhovich, publisher's English record and expanded abstract
for [S1]. \url{https://www.mathnet.ru/eng/im9597}.
\textit{Inspected 27 September 2026; the rational and completed analogues
are distinguished from the discrete integral target.}
\end{itemize}
\end{document}
